% Lösungen Lagebeziehungen (zusammenbau v0.4, Heft)
\einheitenkopf{Lagebeziehungen}

\erg{49}{a) auf der Seite mit größerem Wert – $6 > 3$. \quad b) $2 \cdot 3 - 1 + 4 = 9$, und $9 \neq 7$ – $P$ liegt nicht in $E$. \quad c) Eingesetzt: $2 - 1 + 4 = 5$; wegen $1 < 5 < 7$ liegt $P$ zwischen $E_1$ und $E_2$.}
\erg{50}{a) aus x: $t = 2$, aus z: $s = 3$; y: $0 + 2 + 3 = 5$ – die dritte Gleichung stimmt, $P$ liegt in $E$. \quad b) aus x: $t = 2$, aus y: $s = 2$; z: $0 + 2 + 2 \cdot 2 = 6$, und $6 \neq 5$ – $P$ liegt nicht in $E$. \quad c) $5 \cdot 2 - 2 \cdot 5 = 0$ – $K$ liegt in $W$; $(5 | -2 | 0) \cdot (0 | 0 | 1) = 0$ – die Normalenvektoren stehen senkrecht aufeinander, also steht $W$ senkrecht auf der Fahrbahn. \quad d) $3 + 2 \cdot 4 + 0 = 11$ – $Q$ liegt in $E$; $\vec{PQ} = (2 | 4 | 2) = 2 \cdot (1 | 2 | 1)$ ist ein Vielfaches des Normalenvektors. \quad e) $2 + 3 - 2 = 3$ – $P$ liegt in $E$; $\vec{PQ} = (4 | 2 | -4) = 2 \cdot (2 | 1 | -2)$ ist ein Vielfaches des Normalenvektors, also steht $g$ senkrecht auf $E$. \quad f) „Untersuche, ob die Ebene $W$ die Kante $AE$ schneidet, und gib gegebenenfalls den Punkt an.“ Aus (2): $r = 2$, im Bereich – $W$ schneidet die Kante in $(5 | 0 | 2)$. \quad g) „Untersuche, ob die Ebene $W$ die Kante $BC$ schneidet, und gib gegebenenfalls den Punkt an.“ Aus (2): $r = 1$, im Bereich – $W$ schneidet die Kante in $(2 | 6 | 1)$. \quad h) „Untersuche, ob die Ebene $W$ die Kante $FG$ schneidet, und gib gegebenenfalls den Punkt an.“ Aus (2): $r = 4$, im Bereich – $W$ schneidet die Kante in $(4 | 3 | 2)$. \quad i) Eingesetzt: $1 + 2 \cdot 0,5 + 1 = 3$; wegen $2 < 3 < 6$ liegt $T$ bei $L_1$ auf der vom Ursprung abgewandten Seite und bei $L_2$ auf der Ursprungsseite; alle Koordinaten sind positiv – $T$ liegt im Inneren. \quad j) $S$: ja – $S$ enthält die Ecken $(0 | 0 | 0)$ und $(8 | 4 | 5)$, denn $8 - 2 \cdot 4 = 0$, also die Raumdiagonale durch das Innere. $T$: nein – im Quader ist $x + 2y \ge 0$, gleich null nur auf der Kante mit $x = y = 0$; $T$ berührt das Haus nur dort.}
\erg{51}{a) Punkt gegen Ebene – $P$ wird in die Gleichung von $E_k$ eingesetzt. \quad b) $2k + 2 - 3 = 5$, also $k = 3$. \quad c) z. B. $r = 1$ und $s = 2$, denn $4r - 2s = 0$; $r = s = 0$ ergibt keine Ebenengleichung.}
\erg{52}{a) $A$ und $B$ liefern $6 = 6$, dort fällt $k$ heraus; $C$: $2k = 6$, also $k = 3$. \quad b) z. B. $r = 5$ und $s = -2$, denn $2r + 5s = 0$; $r = s = 0$ ergibt keine Ebenengleichung.}
\erg{53}{a) liegt in der Ebene – $t$ fällt heraus, und $4 = 4$ stimmt. \quad b) $(1 + t) + 2(2 - t) + t = 5$, also $5 = 5$ für jedes $t$ – $g$ liegt in $E$. \quad c) $\vec{u} \cdot \vec{n} = 2 - 2 + 1 = 1 \neq 0$ – genau ein Schnittpunkt.}
\erg{54}{a) $2 + (-1 + 6r) + 2(3 - 3r) = 7$, also $7 = 7$ für jedes $r$ – $s$ liegt in $E$; der Stützpunkt allein hätte nicht gereicht. \quad b) $\vec{u} \cdot \vec{n} = 30 \cdot 1 + 10 \cdot 0 + 1 \cdot (-30) = 0$ – parallel; Stützpunkt: $100 - 30 \cdot 40 = -1100 \neq -800$ – $h$ liegt nicht in $E$, sondern echt parallel. \quad c) nein – die z-Koordinaten $8$ und $6$ hängen nicht von $t$ ab; der Richtungsvektor hat für jedes $t$ die z-Komponente $-2$. \quad d) Die z-Koordinaten von $A$ und $B$ stimmen überein, also hat $\vec{AB} = (2 | -5 | 0)$ die z-Komponente null.}
\erg{55}{a) nein – ein Bereich ist noch zu prüfen: die Koordinaten gegen die Wandmaße. \quad b) Lichtgerade $\vec{x} = (2 | 0 | 4) + t \cdot (1 | 2 | -1)$; $2t = 6$ gibt $t = 3$, Schattenpunkt $(5 | 6 | 1)$; $0 \le 5 \le 6$ und $0 \le 1 \le 3$ – der Schatten liegt auf der Wand. \quad c) Der Schritt schneidet die Lichtgerade durch die Segelecke mit dem Boden; der Schattenpunkt $(5 | 4 | 0)$ hat $x = 5 > 4$ – der Schatten dieser Ecke liegt neben der Terrasse.}
\erg{56}{a) Schritt I bestimmt den Schattenpunkt $(3 | 0 | 2)$ der Ecke $(4 | 2 | 3)$ als Durchstoßpunkt der Lichtgeraden mit der Wandebene; Schritt II zeigt, dass er innerhalb der Wandmaße liegt – der Schatten der Ecke fällt auf die Hauswand. \quad b) Der Schritt schneidet die Lichtgerade durch die Rollo-Ecke mit der Wandebene: $(2 | 0 | 3)$ ist der Schattenpunkt der Ecke; wegen $0 \le 2 \le 5$ und $0 \le 3 \le 4$ liegt er auf der Wand.}
